\section{Research problem and objective}

MRAM compute-in-memory combines nonvolatile weight storage with parallel computation.
Its useful precision and efficiency depend on the circuit that converts cell states and
input activity into a sensed result. A large number of operations per read can reduce
conversion overhead while increasing the demands on current handling, sensing and
interconnect. Existing MRAM prototypes and calibration work demonstrate why the memory
device and readout architecture must be considered together \cite{S08,S23,S24}.

The target of this phase is a \textbf{physical \target\ array}, containing 65,536 MTJ
positions. The original proposal concerns reliability and failure regimes under coupled
numerical precision and hardware nonidealities. The current task is to develop a specific
method or design contribution within that subject, with evidence that exceeds another
parameter sweep. Earlier work is preserved in an archive and supplies no target-array
validation in this draft.

The governing question is:
\begin{quote}
Which prediction, encoding, organization or readout method can improve the cost of useful
reliable computation on the declared target, and what physical mechanism makes the
improvement possible beyond tuned existing methods?
\end{quote}

The question has seven competing method families. Broad techniques such as calibration,
mixed precision, activation grouping and speculative refinement already exist. Their
existence determines the comparison; it does not establish that every narrower question
has been answered. Conversely, using MRAM instead of another resistive technology does
not itself establish originality.

\subsection{Research hierarchy and proposal continuity}

The new research phase has one primary \emph{development} path, not one proven contribution.
Its foundation is matched neural inference and the precision--nonideality failure-regime
analysis specified in the original proposal. F3, a resident row mapping intended to
control decoded cross-cycle error covariance, is the next bounded intervention to test.
F7 investigates an independent Ising field-persistence question only if its physical
mechanism and scope warrant continuation. F1 is a potential evaluation tool; other
families remain recorded alternatives and necessary comparators.

The uploaded 4 November revision states that the original proposal's CA2 milestones
include initial modelling, a baseline network, representative perturbations and
preliminary inference and failure-regime criteria. These are project-planning
deliverables, \emph{not} an independently verified official marking rubric.
No new supervisory endorsement is inferred.

\subsection{Contribution filters and acceptance gates}

Each candidate is screened for (i) useful value at fixed task quality, (ii) a distinct
technical difference, (iii) physical credibility, (iv) bounded FYP scope and (v) advantage
over simpler tuned alternatives after overhead. Three gates precede a large evaluation:
\begin{enumerate}
\item Specify the algorithm, derivation or circuit intervention and its exact difference
from the closest applicable method.
\item Establish one physically justified operating case in the target architecture.
\item Demonstrate useful value in a small comparison against strong baselines, including
calibration, model setup and additional control or hardware cost.
\end{enumerate}
F3 is prioritized for this bounded investigation; no family has passed all three gates or been approved as the final novel contribution. An experiment can reject a defined contribution;
it does not create originality automatically. A failed method remains a rejected method
unless a different, explicitly justified research objective is adopted.

\subsection{Evidence policy}

This draft distinguishes exact software verification, workload structure, illustrative
models, target circuit simulation, validated surrogates, task evaluation and measured
silicon. Operation counts and activity proxies are not energy or reliability measurements.
The reset completed software verification; the 4 November E03M follow-up adds a
constructed mathematical/model example. A separate exploratory 256-wide digits-workload screen has been documented, but it is
not the frozen FYP workload or target circuit evidence. No new validated target circuit
or end-to-end physical inference result is introduced.

Academic sources are used at their actual access depth. A public abstract, inaccessible
submission or reproduced patent description cannot establish the absence of an
implementation feature. The literature register records these boundaries.
